% =============================================================
% SECTION II: RELATED WORK
% =============================================================
\section{Related Work}
\label{sec:related_work}

This section surveys prior work across three intersecting domains: legal AI systems for the Indian context, retrieval-augmented generation, and hardware-rooted trust for document integrity.

\subsection{Legal AI Systems in India}

Recent work in Indian legal informatics has explored conversational assistants, document summarization, and retrieval-based legal guidance. Systems such as LawPal~\cite{panchal2025lawpal}, BharatLex~\cite{bharatlex2025}, and Legal Ally~\cite{borikar2025legalally} demonstrate that retrieval-augmented approaches can improve factual grounding and usability. However, many of these systems remain limited to isolated question answering or document analysis and do not integrate hardware-rooted trust or end-to-end document lifecycle management.

Hybrid reasoning approaches, including the combination of retrieval with structured knowledge representations, have been explored for complex legal queries~\cite{sindhuja2025smartlegal}. Other efforts focus on domain-specific question classification and intent modeling for legal services~\cite{ramesh2025judicial}. These studies underline the need for systems that connect retrieval, analysis, and trustworthy provenance in a unified workflow.

\subsection{Retrieval-Augmented Generation}

RAG, introduced by Lewis \textit{et al.}~\cite{lewis2020rag}, addresses hallucination by conditioning generation on retrieved evidence. Surveys of RAG architectures highlight modular designs that tailor retrieval strategies to query context~\cite{gao2024ragsurvey}. NyayaSahaya follows this modular perspective by separating general legal retrieval from document-specific grounding while keeping implementation details abstract.

\subsection{Hardware-Rooted Trust and Document Integrity}

The integration of physical verification with digital workflows has been studied in the security and privacy literature. Multi-factor authentication and edge-based verification can reduce unauthorized access and improve auditability~\cite{jain2016biometric, shi2016edge}. Standardized signing approaches and integrity ledgers provide strong guarantees of document provenance. NyayaSahaya combines these ideas at a system level to emphasize trust without exposing implementation specifics.
